% © 2026 Ing. David Jaroš — CC BY-NC-ND 4.0
%
% This work is licensed under a Creative Commons Attribution-NonCommercial-NoDerivatives
% 4.0 International License (CC BY-NC-ND 4.0).
%
% See LICENSE.md for full license text.

\documentclass[12pt,a4paper]{article}
\usepackage[a4paper,margin=2.5cm]{geometry}
\usepackage{amsmath,amssymb,tcolorbox,booktabs}
\tcbuselibrary{skins,breakable}
\newtcolorbox{statusbox}[1][]{colback=gray!8!white,colframe=black!60,arc=3pt,boxrule=1pt,title=Status Summary,#1}
\title{Hecke higher-weight check at $p=137$: $k=4,6$}
\author{Ing.\ David Jaroš}
\date{2026-05-12}
\begin{document}
\maketitle

Existing repository values under check:
\[
a_{137}(k=4)=2274,\qquad a_{137}(k=6)=38286.
\]
Derived ratios used by current signal:
\[
R_\mu=\frac{2274}{11}=206.73,\qquad
R_\tau=\frac{38286}{11}=3480.5.
\]
Both satisfy Weil bounds in their weights.

At this stage, this note consolidates the verification protocol and records
that direct independent recomputation (LMFDB/Sage) is still pending in this
sandbox session.

\begin{statusbox}
\textbf{Hlavní výsledek}: Kontrolní rámec pro $k=4,6$ je sjednocen; stávající hodnoty jsou interně konzistentní a Weilově přípustné.\\
\textbf{Klasifikace}: [MC] (independent external recomputation pending).\\
\textbf{Dopad na teorii}: Lepton-mass signal zůstává silný kandidát, ale bez nového externího výpočtu není povýšen na uzavřený důkaz.\\
\textbf{Příští krok}: Provest nezávislý výpočet $a_{137}$ pro příslušné newforms (Sage/LMFDB export) a porovnat bitově.\\
\textbf{Alpha je odvozena}: NE (unconditional)
\end{statusbox}

\end{document}
